\chapter{FX Swaps, Forwards and the Cross-Currency Basis}\label{ch:m2:fx-swaps-forwards-and-the-cross-currency-basis}

At the end of August 2026 a ten-year US Treasury yielded 4.75\% and a ten-year
Japanese government bond 2.94\%. A Japanese insurer that buys the Treasury
earns 1.81 percentage points more, but it owes its policyholders yen, and a fall
of the dollar would wipe that out many times over, so it hedges: it sells the
dollars forward, three months at a time. With the dollar overnight rate at
3.68\% and the yen's under 1\%, each roll of the hedge costs about 2.70
percentage points a year, more than the extra yield, before the market charges
anything extra for lending dollars. The hedged Treasury yields about 2.05\%,
less than the Japanese bond. This chapter is about the instruments that set
that cost: the \omterm{def:m2:fx-swaps-forwards-and-the-cross-currency-basis:forward}{outright forward}, the \omterm{def:m2:fx-swaps-forwards-and-the-cross-currency-basis:swap}{FX swap} that is the most traded
instrument in foreign exchange, the parity that links them to interest rates,
and the basis by which that parity has failed since 2007.

\section{Outright forwards and forward points}

\begin{definition}[Outright forward, forward points]\label{def:m2:fx-swaps-forwards-and-the-cross-currency-basis:forward}
An \emph{outright forward}\index{outright forward} is an agreement to exchange
two currencies at a rate $F$ fixed today, for value on a date after spot. Its
\emph{forward points}\index{forward points} are $F - S$ expressed in pips of the
pair: the market quotes forwards as points added to the spot rate.
\end{definition}

A forward is priced by an argument that uses nothing but deposits. Borrow one
unit of the base currency for the period at its rate $r_b$, sell it spot for $S$
units of the quote currency, and lend those at the quote currency's rate $r_q$:
at maturity you owe $1 + r_b\tau_b$ of the base and own $S(1 + r_q\tau_q)$ of the
quote, where $\tau$ are the accrual fractions in each currency's day count. The
forward that makes the round trip worth nothing is the covered-parity forward.

\begin{definition}[Covered interest parity]\label{def:m2:fx-swaps-forwards-and-the-cross-currency-basis:cip}
\emph{Covered interest parity}\index{covered interest parity} (CIP) is the
relation
\begin{equation}\label{eq:m2:fx-swaps-forwards-and-the-cross-currency-basis:cip}
F \;=\; S\,\frac{1 + r_q\,\tau_q}{1 + r_b\,\tau_b}
\end{equation}
between spot, forward and the two currencies' money-market rates over the same
period.
\end{definition}

The currency with the higher interest rate trades at a forward discount: its
forward buys less of the other than its spot does, and the lower yield of the
other currency is made up by the gain in the exchange rate. USDJPY is quoted in
yen per dollar; with dollar rates above yen rates, its \omterm{def:m2:fx-swaps-forwards-and-the-cross-currency-basis:forward}{forward points} are
negative.

\begin{example}[Three-month USDJPY]\label{ex:m2:fx-swaps-forwards-and-the-cross-currency-basis:fwd}
With USDJPY at 156.87, a dollar rate of 3.68\% (actual/360) and a yen rate of
0.977\% (actual/365), the 91-day forward by \cref{eq:m2:fx-swaps-forwards-and-the-cross-currency-basis:cip}
is 155.8028: \omterm{def:m2:fx-swaps-forwards-and-the-cross-currency-basis:forward}{forward points} of $-106.7$ pips of 0.01 yen. A basis of $-25$ basis
points on the yen leg, as defined below, lowers it to 155.7059, $-116.4$ pips.
The one-day \omterm{def:m2:fx-swaps-forwards-and-the-cross-currency-basis:swap}{tom-next} points are $-1.18$.
\end{example}

\section{FX swaps and short-dated rolls}

\begin{definition}[FX swap, tom-next]\label{def:m2:fx-swaps-forwards-and-the-cross-currency-basis:swap}
An \emph{FX swap}\index{FX swap} is a pair of opposite exchanges of the same
amount of one currency: a near leg, usually at spot, and a far leg at a forward
date, the two rates differing by the \omterm{def:m2:fx-swaps-forwards-and-the-cross-currency-basis:forward}{forward points}. \emph{Tom-next}\index{tom-next}
is the one-day swap from the next business day (tomorrow) to the day after
(next), with which positions are rolled from one spot date to the next.
\end{definition}

An \omterm{def:m2:fx-swaps-forwards-and-the-cross-currency-basis:swap}{FX swap} is a collateralised loan. The party that buys dollars at the near
leg and sells them back at the far leg has borrowed dollars and lent the other
currency for the period, each loan secured by the other; the \omterm{def:m2:fx-swaps-forwards-and-the-cross-currency-basis:forward}{forward points} are
the interest differential. It is by far the largest FX instrument: USD~4
trillion a day in April 2025, 42\% of all FX turnover, mostly for seven days or
less (\cref{ch:m2:the-fx-market}). Banks use it to fund their dollar assets,
investors to hedge their foreign holdings, and a trader who holds a spot
position rolls it every day with a \omterm{def:m2:fx-swaps-forwards-and-the-cross-currency-basis:swap}{tom-next} swap, paying or earning the day's
points: this is the carry of an FX position.

\begin{omfigure}
\centering
\begin{tikzpicture}[font=\small,
  box/.style={draw,thick,rounded corners=2pt,align=center,minimum height=1cm,minimum width=2.8cm,inner sep=3pt},
  arr/.style={-{Stealth[length=2.5mm]},thick}]
\node[box,draw=omDef,fill=omDef!8] (inv) at (0,0) {Japanese investor};
\node[box,draw=omThm,fill=omThm!8] (bk) at (9.2,0) {dollar lender\\(bank)};
\node[font=\footnotesize\bfseries,anchor=west] at (-1.4,0.95) {near leg, spot};
\node[font=\footnotesize\bfseries,anchor=west] at (-1.4,-1.35) {far leg, three months later};
\draw[arr] ([yshift=5pt]inv.east) -- node[above,font=\footnotesize] {JPY~1\,568.70 million} ([yshift=5pt]bk.west);
\draw[arr] ([yshift=-5pt]bk.west) -- node[below,font=\footnotesize] {USD~10 million} ([yshift=-5pt]inv.east);
\draw[arr,dashed] ([yshift=-1.8cm]bk.west) -- node[above,font=\footnotesize] {JPY~1\,557.06 million} ([yshift=-1.8cm]inv.east);
\draw[arr,dashed] ([yshift=-2.3cm]inv.east) -- node[below,font=\footnotesize] {USD~10 million} ([yshift=-2.3cm]bk.west);
\end{tikzpicture}

\omcaption{A three-month USDJPY \omterm{def:m2:fx-swaps-forwards-and-the-cross-currency-basis:swap}{FX swap} at the forward of
\cref{ex:m2:fx-swaps-forwards-and-the-cross-currency-basis:fwd} with a $-25$
basis-point basis. The investor borrows dollars and lends yen, each loan secured
by the other; it gets back JPY~11.64 million fewer than it paid, the price of
the dollars. Illustrative.}\label{fig:m2:fx-swaps-forwards-and-the-cross-currency-basis:legs}
\end{omfigure}

\section{Covered interest parity}

Before 2007 CIP held within transaction costs, because any gap was a riskless
profit for a bank: if dollars were cheaper to borrow through \omterm{def:m2:fx-swaps-forwards-and-the-cross-currency-basis:swap}{FX swaps} than in
the money market, a bank would borrow them through swaps and lend them in cash,
and its trades would close the gap. Economists at the BIS called it ``the
closest thing to a physical law in international finance''. The same
economists documented, in 2016, that it had been systematically violated since
the financial crisis of 2007--2008.

\begin{proposition}[Forward points by tenor]\label{prop:m2:fx-swaps-forwards-and-the-cross-currency-basis:points}
Under CIP with flat rates, the \omterm{def:m2:fx-swaps-forwards-and-the-cross-currency-basis:forward}{forward points} of a pair grow almost in
proportion to the tenor, $F - S \approx S\,(r_q - r_b)\,\tau$ for short tenors,
and a basis $b$ on the quote leg adds approximately $S\,b\,\tau$.
\end{proposition}

\begin{proof}
Expand \cref{eq:m2:fx-swaps-forwards-and-the-cross-currency-basis:cip} to first
order in $\tau$: $F/S \approx 1 + r_q\tau - r_b\tau$. With $r_q + b$ in place of
$r_q$ the term $b\tau$ appears.
\end{proof}

\begin{omfigure}
\begin{tikzpicture}
\begin{axis}[width=11cm,height=4.8cm,
  xlabel={days to the far leg},ylabel={USDJPY forward points},
  tick label style={font=\small},label style={font=\small},
  xmin=0,xmax=370,ymin=-480,ymax=10,grid=major,grid style={gray!25},
  legend columns=2,legend style={font=\footnotesize,at={(0.5,-0.32)},anchor=north,draw=none,fill=none,/tikz/every even column/.append style={column sep=8pt}},
  table/col sep=comma]
\addplot[omDef,very thick,mark=*,mark size=1.3pt] table[x=days,y=nobasis]{figdata/markets-2/16-fx-swaps-forwards-and-the-cross-currency-basis/points.csv};
\addplot[omThm,very thick,dashed,mark=*,mark size=1.3pt] table[x=days,y=basis]{figdata/markets-2/16-fx-swaps-forwards-and-the-cross-currency-basis/points.csv};
\legend{covered interest parity,with a $-25$ bp basis on the yen leg}
\end{axis}
\end{tikzpicture}

\omcaption{USDJPY \omterm{def:m2:fx-swaps-forwards-and-the-cross-currency-basis:forward}{forward points}, in pips of 0.01 yen, from one week to one year,
with flat money-market rates of 3.68\% in dollars and 0.977\% in yen, at a spot of
156.87. The points grow with the tenor; a negative basis on the yen makes them
more negative, the dollar dearer forward. Illustrative; data: the chapter's
tutorial.}\label{fig:m2:fx-swaps-forwards-and-the-cross-currency-basis:points}
\end{omfigure}

\section{The cross-currency basis}

\begin{definition}[Cross-currency basis swap, cross-currency basis]\label{def:m2:fx-swaps-forwards-and-the-cross-currency-basis:basis}
A \emph{cross-currency basis swap}\index{cross-currency basis swap} exchanges
notionals in two currencies at the start at the spot rate, exchanges floating
interest in each currency during its life, one leg with a spread, and
re-exchanges the same notionals at maturity. The spread is the
\emph{cross-currency basis}\index{cross-currency basis} $b$, quoted on the
non-dollar leg; for short tenors it is measured from \omterm{def:m2:fx-swaps-forwards-and-the-cross-currency-basis:swap}{FX swaps} as the gap between
the yen (or euro) rate implied by \omterm{def:m2:fx-swaps-forwards-and-the-cross-currency-basis:forward}{forward points} and the cash rate:
\[
b \;=\; \left(\frac{F}{S}\,(1 + r_b\tau_b) - 1\right)\frac{1}{\tau_q} - r_q.
\]
\end{definition}

A negative basis on the yen means that lending yen and borrowing dollars through
swaps earns less than the yen money-market rate: dollars are dearer through
swaps than in cash. Since 2007 the basis against the dollar has been negative for
most currencies, notably the euro and the yen, and positive for a few, such as
the Australian dollar. The BIS explanation combines demand and limits. Demand:
investors such as Japanese insurers hedge large dollar holdings by borrowing
dollars in swaps, and companies that issue bonds in euros and swap the proceeds
into dollars add to it. Limits: the banks that would arbitrage the gap now pay
for the balance sheet the trade uses, in capital and leverage ratios, and do so
only for a spread. Since 2014 the basis has jumped at quarter ends, with \omterm{def:m2:repo-and-specials:rate}{repo
rates}, when banks shrink their balance sheets for their reports.

For the investor of the hook the basis is a direct cost. Rolling a three-month
hedge of a dollar bond costs, a year, approximately the dollar rate minus the
yen rate plus the yen basis taken with its sign reversed:
\begin{equation}\label{eq:m2:fx-swaps-forwards-and-the-cross-currency-basis:hedged}
y^{\mathrm{hedged}} \;\approx\; y_{\mathrm{USD}} - \bigl(r_{\mathrm{USD}} - r_{\mathrm{JPY}} - b\bigr).
\end{equation}

\begin{omfigure}
\begin{tikzpicture}
\begin{axis}[width=11cm,height=5cm,
  ylabel={percent},
  tick label style={font=\small},label style={font=\small},
  xmin=2018.2,xmax=2026.75,ymin=-2,ymax=5.5,grid=major,grid style={gray!25},
  xtick={2019,2020,2021,2022,2023,2024,2025,2026},xticklabel style={/pgf/number format/1000 sep={}},
  legend columns=3,legend style={font=\footnotesize,at={(0.5,-0.18)},anchor=north,draw=none,fill=none,/tikz/every even column/.append style={column sep=8pt}},
  table/col sep=comma]
\addplot[omDef,thick] table[x=t,y=ust]{figdata/markets-2/16-fx-swaps-forwards-and-the-cross-currency-basis/hedged.csv};
\addplot[omThm,very thick] table[x=t,y=hedged]{figdata/markets-2/16-fx-swaps-forwards-and-the-cross-currency-basis/hedged.csv};
\addplot[omProp,thick,dashed] table[x=t,y=jgb]{figdata/markets-2/16-fx-swaps-forwards-and-the-cross-currency-basis/hedged.csv};
\draw[gray] (axis cs:2018.2,0) -- (axis cs:2026.75,0);
\legend{10-year Treasury,hedged into yen (zero basis),10-year JGB}
\end{axis}
\end{tikzpicture}

\omcaption{A ten-year Treasury for a yen investor, April 2018 to August 2026:
unhedged, hedged by rolling short-dated \omterm{def:m2:fx-swaps-forwards-and-the-cross-currency-basis:swap}{FX swaps} with no basis (the Treasury
yield less SOFR plus the Japanese call rate), and the ten-year JGB. Any negative
basis lowers the hedged line further. Month-end Treasury yields and SOFR, monthly
averages for Japan. Data: FRED series DGS10, SOFR, IRSTCI01JPM156N and
IRLTLT01JPM156N.}\label{fig:m2:fx-swaps-forwards-and-the-cross-currency-basis:hedged}
\end{omfigure}

\begin{dated}{2026-09}{The inputs of the hook}\label{dat:m2:fx-swaps-forwards-and-the-cross-currency-basis:inputs}
Ten-year Treasury 4.75\% and SOFR 3.68\% on 31 August 2026; Japanese call-money
rate 0.977\% and ten-year JGB 2.94\% as August 2026 averages; USDJPY 156.87 on 18
September 2026 (Federal Reserve H.10). Overnight rates stand in for three-month
rates, and no official series of the basis is used: its levels in this chapter
are illustrative.
\end{dated}

\Cref{fig:m2:fx-swaps-forwards-and-the-cross-currency-basis:hedged} shows the
consequence. When the Federal Reserve's rates were far above Japan's, from late
2018 to early 2020 and again from November 2022, the hedged Treasury yielded less
than the Japanese bond even before any basis: in 57 of the 101 months shown. The investor
then chooses between a lower hedged yield, an unhedged position that bets on the
yen, or staying at home.

\section{Dollar funding stress}

Non-US banks hold large dollar assets funded partly through \omterm{def:m2:fx-swaps-forwards-and-the-cross-currency-basis:swap}{FX swaps}. When
dollar funding dries up, as in 2008 and March 2020, the basis widens sharply,
and the central bank that issues the dollar lends it to other central banks,
which lend it on to their banks. The Federal Reserve first set up such swap
lines on 12 December 2007, with the ECB for up to USD~20 billion and the Swiss
National Bank for USD~4 billion; on 13 October 2008 the lines with the Bank of
England, the ECB and the SNB were increased ``to accommodate whatever quantity of
U.S. dollar funding is demanded''. In March 2020 the standing lines with five
central banks were cut in price to the dollar OIS rate plus 25 basis points on 15
March, and temporary lines with nine more were opened on 19 March
(\cref{fig:m2:fx-swaps-forwards-and-the-cross-currency-basis:swaplines}).

\begin{omfigure}
\begin{tikzpicture}
\begin{axis}[width=11cm,height=4.6cm,
  ylabel={USD billion},
  tick label style={font=\small},label style={font=\small},
  xmin=2007,xmax=2026.75,ymin=0,ymax=650,grid=major,grid style={gray!25},
  xtick={2008,2010,2012,2014,2016,2018,2020,2022,2024,2026},xticklabel style={/pgf/number format/1000 sep={}},
  table/col sep=comma]
\addplot[omThm,very thick,const plot] table[x=t,y=bn]{figdata/markets-2/16-fx-swaps-forwards-and-the-cross-currency-basis/swaplines.csv};
\end{axis}
\end{tikzpicture}

\omcaption{Dollars lent by the Federal Reserve to foreign central banks through
swap lines, month ends, 2007 to August 2026. The peaks are the dollar funding
crises of late 2008 (a weekly high of USD~583 billion on 17 December 2008) and
spring 2020 (USD~449 billion on 27 May 2020). Data: FRED series SWPT, from the H.4.1
release.}\label{fig:m2:fx-swaps-forwards-and-the-cross-currency-basis:swaplines}
\end{omfigure}

The swap lines cap the basis in a crisis, since a bank with access to its central
bank's dollar auctions will not pay much more than their rate. They do not remove
it in normal times, when the drivers are hedging demand and the price of balance
sheet.

\section{Tutorial: forward points and the implied basis}\label{tut:m2:fx-swaps-forwards-and-the-cross-currency-basis:basis}

\begin{tutorial}
\textbf{Goal.} Compute \omterm{def:m2:fx-swaps-forwards-and-the-cross-currency-basis:forward}{forward points} from two money-market rates, back out the
basis implied by market points, and compute a hedged yield.
\textbf{End state:}
\cref{ex:m2:fx-swaps-forwards-and-the-cross-currency-basis:fwd},
\cref{fig:m2:fx-swaps-forwards-and-the-cross-currency-basis:points,fig:m2:fx-swaps-forwards-and-the-cross-currency-basis:hedged}
and the numbers of the weekend problem.
\begin{tutsteps}
\item \textbf{The forward and the basis}, by
\cref{eq:m2:fx-swaps-forwards-and-the-cross-currency-basis:cip} and its inverse.
\omcode{code/firm/fxfwd/firm_fxfwd.py}{10}{30}{Forward by covered parity, forward points, implied rate and basis.}\label{lst:m2:fx-swaps-forwards-and-the-cross-currency-basis:fwd}
\item \textbf{The hedge}, by
\cref{eq:m2:fx-swaps-forwards-and-the-cross-currency-basis:hedged}.
\omcode{code/firm/fxfwd/firm_fxfwd.py}{39}{47}{Hedge cost and hedged yield.}\label{lst:m2:fx-swaps-forwards-and-the-cross-currency-basis:hedge}
\item \textbf{Run} \texttt{fxfwd\_demo.three\_month()},
\texttt{fxfwd\_demo.hedged\_now()} and \texttt{fig\_fxfwd.py}.
\end{tutsteps}
\textbf{What to change next.} Replace the overnight rates by term rates of the
hedge's tenor, and compare the hedged yield from rolling three-month swaps with
that from a one-year forward.
\end{tutorial}

\section{Build: the forward and basis calculator}\label{bld:m2:fx-swaps-forwards-and-the-cross-currency-basis:fxfwd}

\begin{build}
\bfield{Purpose} The miniature firm quotes forwards and swaps to clients, funds
dollar positions through swaps, and hedges investors' currency exposure: it needs
forwards from curves, the basis from market points, and hedged yields.
\bfield{Interface} \texttt{forward(spot, r\_quote, r\_base, days, basis\_quote,
basis\_base, dc\_quote, dc\_base)}; \texttt{points(fwd, spot, pip)};
\texttt{implied\_rate\_quote}; \texttt{basis\_on\_quote}; \texttt{fx\_swap\_legs};
\texttt{hedge\_cost}; \texttt{hedged\_yield}.
\bfield{Rules} Simple money-market rates with each currency's day count; basis
on the non-dollar leg; \omterm{def:m2:fx-swaps-forwards-and-the-cross-currency-basis:forward}{forward points} in the pair's pips.
\bfield{Acceptance tests} \texttt{code/firm/fxfwd/tests/}: CIP forward and its
inverse; the basis recovered from a forward built with it; \omterm{def:m2:fx-swaps-forwards-and-the-cross-currency-basis:swap}{FX swap} legs that
net to the interest differential; the hedged yield lower with a negative basis.
\bfield{Stretch} Forward curves from full discount curves in each currency
(\cref{bld:m2:interest-rate-swaps:curve}); broken dates; a \omterm{def:m2:fx-swaps-forwards-and-the-cross-currency-basis:basis}{cross-currency basis
swap} priced on two curves with the basis as a spread.
\end{build}

\begin{omsources}
\item C.~Borio, R.~McCauley, P.~McGuire and V.~Sushko, ``Covered interest parity
lost: understanding the cross-currency basis'', \emph{BIS Quarterly Review},
September 2016.
\item Bank for International Settlements, Triennial Central Bank Survey 2025.
\item Federal Reserve, press releases of 12 December 2007, 13 October 2008, 15
and 19 March 2020; H.4.1 and H.10 releases.
\item FRED series DGS10, SOFR, IRSTCI01JPM156N, IRLTLT01JPM156N and SWPT.
\end{omsources}

\section{Exercises}

\begin{exercise}[$\star$]\label{exo:m2:fx-swaps-forwards-and-the-cross-currency-basis:1}
EURUSD is 1.1464, the euro rate 2.00\% and the dollar rate 3.68\%, both
actual/360. Give the 90-day forward and its points.
\end{exercise}

\begin{exercise}[$\star$]\label{exo:m2:fx-swaps-forwards-and-the-cross-currency-basis:2}
Why does the currency with the higher interest rate trade at a forward discount?
\end{exercise}

\begin{exercise}[$\star$]\label{exo:m2:fx-swaps-forwards-and-the-cross-currency-basis:3}
In the swap of \cref{fig:m2:fx-swaps-forwards-and-the-cross-currency-basis:legs},
who has borrowed which currency, and what is the yen cost of the dollars?
\end{exercise}

\begin{exercise}[$\star\star$]\label{exo:m2:fx-swaps-forwards-and-the-cross-currency-basis:4}
The market's three-month USDJPY forward is 155.7059 with the rates of
\cref{ex:m2:fx-swaps-forwards-and-the-cross-currency-basis:fwd}. What basis does
it imply, and what does its sign mean?
\end{exercise}

\begin{exercise}[$\star\star$]\label{exo:m2:fx-swaps-forwards-and-the-cross-currency-basis:5}
A trader long USD~100 million against yen rolls the position with \omterm{def:m2:fx-swaps-forwards-and-the-cross-currency-basis:swap}{tom-next}.
Using the \omterm{def:m2:fx-swaps-forwards-and-the-cross-currency-basis:swap}{tom-next} points of \cref{ex:m2:fx-swaps-forwards-and-the-cross-currency-basis:fwd},
what does one roll earn or cost, in yen?
\end{exercise}

\begin{exercise}[$\star\star$]\label{exo:m2:fx-swaps-forwards-and-the-cross-currency-basis:6}
Why do quarter ends move the basis, and why did the swap lines of March 2020
narrow it?
\end{exercise}

\begin{exercise}[$\star\star\star$]\label{exo:m2:fx-swaps-forwards-and-the-cross-currency-basis:7}
\emph{Coding.} With \texttt{hedged\_series}, find the first month in the data
when the zero-basis hedged Treasury yielded less than the JGB, and count such
months.
\end{exercise}

\begin{exercise}[$\star\star\star$]\label{exo:m2:fx-swaps-forwards-and-the-cross-currency-basis:8}
\emph{Find the flaw.} ``The basis is negative, so a bank can borrow dollars in
the money market, lend them through \omterm{def:m2:fx-swaps-forwards-and-the-cross-currency-basis:swap}{FX swaps} and lock in a riskless profit;
this will close the basis within days.'' Correct it.
\end{exercise}

\section{Problem: The Hedged Yield}

\begin{problem}[{Weekend problem --- what a Treasury yields to a yen investor}]\label{pb:m2:fx-swaps-forwards-and-the-cross-currency-basis:1}
A Japanese life insurer considers buying ten-year Treasuries at 4.75\% on 31
August 2026 and hedging them into yen with rolling three-month \omterm{def:m2:fx-swaps-forwards-and-the-cross-currency-basis:swap}{FX swaps}. Dollar
rates are 3.68\%, yen rates 0.977\%, USDJPY 156.87; the ten-year JGB yields 2.94\%.

\medskip\noindent\textbf{Part I --- The hedge.}
\begin{enumerate}
\item Give the three-month forward with no basis, and its points.
\item Give the annualised cost of the hedge with no basis.
\item Give the hedged yield with no basis.
\item Give it with a basis of $-25$ and $-50$ basis points.
\item Compare each with the JGB.
\end{enumerate}

\medskip\noindent\textbf{Part II --- The mechanics.}
\begin{enumerate}[resume]
\item What does the insurer do at each roll?
\item What happens to the hedge's cash flows if the dollar falls 10\% in a quarter?
\item Why does the insurer need yen liquidity, and when?
\item How does a one-year forward hedge differ from rolling three-month swaps?
\item What would make the hedge cheaper?
\end{enumerate}

\medskip\noindent\textbf{Part III --- The alternatives.}
\begin{enumerate}[resume]
\item What does the insurer earn unhedged, and what does it risk?
\item What break-even does the dollar need, over a year, for the unhedged position to beat the JGB?
\item Why might it hedge only half?
\item Why do such flows widen the basis?
\item Who takes the other side of its \omterm{def:m2:fx-swaps-forwards-and-the-cross-currency-basis:swap}{FX swaps}?
\end{enumerate}

\medskip\noindent\textbf{Part IV --- Judgement.}
\begin{enumerate}[resume]
\item When did hedged Treasuries last beat JGBs in the data, and why then?
\item How would a cut of 100 basis points by the Federal Reserve change the answer?
\item Why is the basis the insurer pays not a free lunch for the bank that lends the dollars?
\item State the \emph{named result}: the hedged yield of the Treasury for the insurer.
\item In one sentence: what does a hedged investor in foreign bonds earn?
\end{enumerate}
\end{problem}

\section{Interview questions}

\begin{interviewq}[$\star$ \iqroles{trader, researcher}]\label{iq:m2:fx-swaps-forwards-and-the-cross-currency-basis:1}
Derive the forward exchange rate from spot and interest rates.
\end{interviewq}

\begin{interviewq}[$\star$ \iqroles{bank, trader}]\label{iq:m2:fx-swaps-forwards-and-the-cross-currency-basis:2}
What is an \omterm{def:m2:fx-swaps-forwards-and-the-cross-currency-basis:swap}{FX swap}, and why is it the most traded FX instrument?
\end{interviewq}

\begin{interviewq}[$\star\star$ \iqroles{researcher, trader}]\label{iq:m2:fx-swaps-forwards-and-the-cross-currency-basis:3}
Why has \omterm{def:m2:fx-swaps-forwards-and-the-cross-currency-basis:cip}{covered interest parity} failed since 2008?
\end{interviewq}

\begin{interviewq}[$\star\star$ \iqroles{trader}]\label{iq:m2:fx-swaps-forwards-and-the-cross-currency-basis:4}
How would you hedge a portfolio of euro bonds for a dollar investor, and what
does the hedge earn or cost?
\end{interviewq}

\begin{interviewq}[$\star\star$ \iqroles{bank}]\label{iq:m2:fx-swaps-forwards-and-the-cross-currency-basis:5}
What happens in the \omterm{def:m2:fx-swaps-forwards-and-the-cross-currency-basis:swap}{FX swap} market at a quarter end, and what would a bank do
about it?
\end{interviewq}

\begin{interviewq}[$\star\star\star$ \iqroles{developer, researcher}]\label{iq:m2:fx-swaps-forwards-and-the-cross-currency-basis:6}
Design a service that publishes \omterm{def:m2:fx-swaps-forwards-and-the-cross-currency-basis:forward}{forward points} for fifty pairs and all standard
tenors, and the implied basis, from rate curves and market points.
\end{interviewq}
